\chapter{建议阅读路径 Suggested Reading Paths}
\section*{线性课程路径}
第一次系统学习 advanced LLM agents 时，建议先按 Chapter 1--13 顺序阅读 Spring 2025 主线，再继续读 Chapter 14--16 的扩展课程章节。这样能先建立统一 vocabulary，再逐步进入 code、multimodality、formal mathematics 与 security，最后再把这些内容接到 Berkeley/Stanford 2024--2025 的更大课程网络中。
\section*{Reasoning-first 路径}
若重点是 reasoning、judge、search 与 theorem proving 的共通结构，建议读 Chapter 1, 2, 3, 4, 9, 10, 11, 13，再回看 Chapter 15 和 Chapter 16 中关于 evaluation 与 self-improvement 的扩展部分。
\section*{Agent systems builder 路径}
若重点是 workflow、tool use、execution-based evaluation 与部署，建议读 Chapter 1, 5, 6, 7, 8, 13，再补 Chapter 14 和 Chapter 15。
\section*{Formal methods 路径}
若来自 theorem proving / verification 背景，建议先读 Chapter 9, 10, 11, 12, 13，再回读 Chapter 2 和 Chapter 6。
\section*{安全与安全防护路径}
若最关心 prompt injection、memory poisoning、privilege control 与 trustworthy deployment，建议读 Chapter 1, 6, 8, 12, 13, 15，并配合 Glossary、Notation 与 Omission Log 一起看。
\section*{课程扩展路径}
若你已经读完 Spring 2025 主线，想系统补 Berkeley/Stanford 相邻课程，请连续阅读 Chapter 14, 15, 16；再回查对应 supplement workspaces 中的 source manifests 和 omission logs，看哪些材料来自官方课程页、哪些来自官方/public playlist、哪些只能以 instructor-affiliated 但非官方的形式记录。
\section*{如何回查 provenance}
每章的 recording、slides、readings、figures 与 evaluator scores 都在 \texttt{textbook\_source\_manifest.json} 和 lecture/supplement workspace sidecars 中可追溯。正文遇到术语边界不确定时，优先回查 Glossary、Notation 与 Figure Provenance。
